% ================================================================================================ apéndices
\chapter{Equipos del OAN}\label{ap:oan}

Estos son los equipos que trae \boton{Dispositivos}. Los datos de los telescopios son los nominales del fabricante;
comprueba la focal de tu tubo en su etiqueta.

\begin{table}[H]
\centering\small
\begin{tabularx}{\linewidth}{@{}l r r Y@{}}
\toprule
Telescopio & Apertura & Focal & Notas\\
\midrule
Celestron NexStar 8SE & 203~mm & 2032~mm (f/10) & Schmidt-Cassegrain; el de las mediciones del grupo.\\
Meade 10″ (LX200) & 254~mm & 2500~mm (f/10) & Con un reductor f/6,3 la focal es de unos 1600~mm.\\
\bottomrule
\end{tabularx}
\caption{Telescopios.}
\end{table}

\begin{table}[H]
\centering\small
\begin{tabularx}{\linewidth}{@{}l l r Y@{}}
\toprule
Cámara & Sensor & Píxel & En video\\
\midrule
Nikon D5600 & APS-C 23,5~×~15,6~mm, 6000~×~4000 & 3,92~µm & Reduce el sensor a 1920~px: píxel efectivo de 12,24~µm.\\
Canon T5i (700D) & APS-C 22,3~×~14,9~mm, 5184~×~3456 & 4,30~µm & Reduce el sensor a 1920~px: píxel efectivo de 11,61~µm.\\
QHY174M-GPS & IMX174, 11,25~×~7,03~mm, 1920~×~1200 & 5,86~µm & Graba con los píxeles reales.\\
\bottomrule
\end{tabularx}
\caption{Cámaras.}
\end{table}

\begin{table}[H]
\centering\small
\begin{tabular}{@{}l c c c@{}}
\toprule
 & Nikon D5600 (video 1920) & Canon T5i (video 1920) & QHY174M-GPS\\
\midrule
Celestron NexStar 8SE & 1,242 & 1,179 & 0,595\\
Meade 10″ (2500~mm) & 1,010 & 0,958 & 0,483\\
\bottomrule
\end{tabular}
\caption{Escalas de placa resultantes, en $''/\text{px}$.}
\end{table}

\textbf{Máscaras.} La del grupo, para el 8SE, tiene orificios de $D = 4{,}2$~cm separados $B = 14{,}4$~cm. La del
protocolo «Determina la nitidez del cielo nocturno» de ANDEAN ROAD, para telescopios de 8 pulgadas, $D = 2{,}5$~cm y
$B = 14$~cm. \emph{Conservar D y B} deja los valores que ya escribiste en \menu{Parámetros}.

% ------------------------------------------------------------------------------------------------ versiones
\chapter{Historial de versiones}\label{ap:versiones}

\begin{table}[H]
\centering\small
\begin{tabularx}{\linewidth}{@{}l Y@{}}
\toprule
Versión & Cambios principales\\
\midrule
1.0.0 & Primera versión: extracción de cuadros, medición vectorizada, seeing con el modelo de Tokovinin, modo educativo,
historial.\\
1.1 & Animación de los centroides; Fredy Esteban Cruz Ospina como coautor; licencia PolyForm Noncommercial.\\
1.2 & Tres idiomas; los mismos controles de reproducción en toda la aplicación; ejecutables para descomprimir y abrir;
historial y perfiles compartidos entre modos; exportación con nombres automáticos.\\
1.3 & Detector de dos grupos (puntos grandes y en anillo); seguimiento cuadro a cuadro y tamaño automático; no se procesa
sin escala de placa; calculadora con el píxel efectivo del video; unidades; perfiles editables; modos Imágenes y Satélite;
animaciones del modo educativo; Saray Juliana Pinto Sanabria como coautora.\\
1.4 & Calibración guiada por las cabeceras, para imágenes sueltas y carpetas completas; botón de los equipos del OAN;
interfaz al 100\,\% por defecto; primer manual completo.\\
\textbf{1.5.0} & \textbf{Medición:} el modo se llama \emph{Medición del seeing}; equipos propios (telescopio, cámara y
máscara) guardados junto a los del OAN; rangos permitidos a la vista; gráficas interactivas (valores con el cursor, ver en
grande, guardar y copiar); gráficas vacías cuando no se detecta la estrella; guardar en el historial varias veces;
exportar en dos pasos con \emph{Generar}; aviso antes de perder un análisis sin guardar; la explicación de las
correlaciones pasa al modo educativo.
\textbf{Satélite:} una noche completa cada 30 o 60 minutos; seeing con los canales 9, 13 y 16 y criterio del Sol; lectura
solo de la zona; mapa sobre el recorte; una subcarpeta por evaluación; la página \emph{Cómo funciona}.
\textbf{Imágenes:} regiones circulares y elípticas; escala variable como en DS9; FWHM y seeing de una estrella; gráficas
de la imagen; diagnóstico del flat; comparar antes y después.
\textbf{Manual} nuevo, con la teoría, capturas numeradas y el código comentado.\\
1.5.1 & José Alejandro López Hernández como autor; los cuatro autores aparecen de la misma forma.\\
\bottomrule
\end{tabularx}
\caption{Versiones de TurbuLENS.}
\end{table}

% ------------------------------------------------------------------------------------------------ créditos
\chapter{Créditos, licencia y cómo citar}\label{ap:creditos}

\section*{Autores}

\textbf{Diego Fernando Cagua Lozada}, \textbf{Fredy Esteban Cruz Ospina}, \textbf{Saray Juliana Pinto Sanabria} y
\textbf{José Alejandro López Hernández}, estudiantes del pregrado de Física, Departamento de Física, Universidad Nacional de Colombia, sede
Bogotá. TurbuLENS se desarrolló para la asignatura Técnicas Observacionales de la Maestría en Astronomía del Observatorio
Astronómico Nacional.

\section*{Desarrollo con inteligencia artificial}

La interfaz y las funcionalidades de TurbuLENS se desarrollaron con ayuda de Claude (Anthropic), un asistente de
inteligencia artificial, bajo la dirección y la revisión de los autores, que definieron el método, validaron los
resultados con datos reales y sintéticos y son responsables del programa.

\section*{Datos y componentes de terceros}

Imágenes del satélite: NOAA GOES-R Series, producto ABI L2 \emph{Cloud and Moisture Imagery} (CMIPF), del repositorio
NOAA Open Data Dissemination en Amazon Web Services. Mapas: Natural Earth (dominio público). TurbuLENS usa Qt para Python
(PySide6, LGPL), FFmpeg a través de PyAV, NumPy, Pillow, openpyxl y pyfive (BSD); cada componente conserva su licencia,
como se detalla en \ruta{THIRD\_PARTY\_NOTICES.md}.

\section*{Licencia}\label{sec:licencia}

PolyForm Noncommercial 1.0.0. Puedes usar, estudiar, modificar y compartir TurbuLENS con fines no comerciales (estudio,
investigación, enseñanza, divulgación), conservando la licencia y el crédito a los autores. El uso comercial requiere
permiso escrito de los autores.

\section*{Cómo citar}\label{sec:citar}

Si usas TurbuLENS en un trabajo, una tesis o una publicación, cítalo así:
\begin{tcolorbox}[enhanced,colback=suave,colframe=suave,arc=2mm,left=4mm,right=4mm,fontupper=\small]
Cagua Lozada, D. F., Cruz Ospina, F. E., Pinto Sanabria, S. J., \& López Hernández, J. A. (2026). \emph{TurbuLENS} (versión 1.5.1) [Software].
Observatorio Astronómico Nacional, Universidad Nacional de Colombia. \url{https://github.com/DiegoCagua/TurbuLENS}
\end{tcolorbox}
El repositorio incluye un archivo \ruta{CITATION.cff} con estos datos, que GitHub muestra como «Cite this repository».

\captura{acerca.png}{\menu{Acerca de}: idioma, tamaño de la interfaz, créditos, licencia, cómo citar y referencias.}

% ------------------------------------------------------------------------------------------------ referencias
\chapter{Referencias}\label{ap:referencias}

\begin{enumerate}[label={[\arabic*]},leftmargin=*,itemsep=0.5em]
\item A. Tokovinin, «From differential image motion to seeing», \emph{Publications of the Astronomical Society of the
Pacific} \textbf{114}, 1156–1166 (2002). \url{https://arxiv.org/abs/astro-ph/0111022}
\item M. Sarazin y F. Roddier, «The ESO differential image motion monitor», \emph{Astronomy \& Astrophysics}
\textbf{227}, 294–300 (1990).
\item D. L. Fried, «Optical resolution through a randomly inhomogeneous medium for very long and very short exposures»,
\emph{Journal of the Optical Society of America} \textbf{56}, 1372–1379 (1966).
\item A. N. Kolmogorov, «The local structure of turbulence in incompressible viscous fluid for very large Reynolds
numbers», \emph{Doklady Akademii Nauk SSSR} \textbf{30}, 301 (1941).
\item F. Roddier, «The effects of atmospheric turbulence in optical astronomy», \emph{Progress in Optics} \textbf{19},
281–376 (1981).
\item S. Cavazzani, S. Ortolani y V. Zitelli, «Site testing at astronomical sites: seeing evaluation from satellite based
data» (2012). \url{https://arxiv.org/abs/1112.3497}
\item NOAA/NASA, \emph{GOES-R Series Product Definition and Users' Guide}, vol. 3 (Level 1b) y vol. 5 (Level 2+).
\item NOAA, aviso de la transición de GOES-19 a GOES-Este, 4 de abril de 2025.
\item ANDEAN ROAD, Oficina Regional Andina de la Astronomía para el Desarrollo, \emph{Determina la nitidez del cielo
nocturno en tu lugar de observación} (protocolo, 2026).
\item IRAF, documentación de la tarea \texttt{display}: el algoritmo ZSCALE (el que usa SAOImage DS9).
\item K. Levenberg, «A method for the solution of certain non-linear problems in least squares», \emph{Quarterly of
Applied Mathematics} \textbf{2}, 164–168 (1944); D. W. Marquardt, \emph{Journal of the Society for Industrial and Applied
Mathematics} \textbf{11}, 431–441 (1963).
\item The Astronomical Almanac, «Low precision formulas for the Sun», sección C.
\end{enumerate}
